% Generated draft: requires mathematical and source review.
\documentclass[11pt]{article}
\usepackage[margin=25mm]{geometry}
\usepackage{fontspec}
\setmainfont{Latin Modern Roman}
\usepackage{amsmath,amssymb,mathtools,unicode-math}
\setmathfont{Latin Modern Math}
\usepackage{enumitem,needspace,xurl,hyperref}
\setlength{\parindent}{0pt}
\setlength{\parskip}{5pt}
\setlength{\emergencystretch}{3em}
\newcommand{\N}{\mathbb{N}}
\newcommand{\Z}{\mathbb{Z}}
\newcommand{\Q}{\mathbb{Q}}
\newcommand{\R}{\mathbb{R}}
\newcommand{\C}{\mathbb{C}}
\newcommand{\F}{\mathbb{F}}
\newcommand{\A}{\mathbb{A}}
\newcommand{\PP}{\mathbb{P}}
\newcommand{\E}{\mathbb{E}}
\newcommand{\eps}{\varepsilon}
\newcommand{\dd}{\,\mathrm{d}}
\begin{document}
\section*{UnsolvedMath ID: 20000823}
\textbf{Source problem code:} catalogue entry
\subsection{Definitions and mathematical statement}
Fix integers \(d\ge1\), \(g\), and \(n\ge2\). Write \(\operatorname{Hilb}_{d,g}(\PP^n)\) for the Hilbert scheme with polynomial \(dt+1-g\), so \(d\) is degree and \(g\) arithmetic genus. The source asks for a good, preferably sharp, lower bound \(L(d,g,n)\) such that every component in its intended class has dimension at least \(L(d,g,n)\). The source does not specify whether all subschemes with this polynomial are allowed or only components whose general member is a smooth connected curve, nor a ground field. For the usual smooth-curve deformation interpretation, work over an algebraically closed characteristic-zero field and keep this interpretation explicit.

For a smooth connected curve \(C\subset\PP^n\), the normal Euler characteristic is
\[
 E_n(d,g)=(n+1)d+(n-3)(1-g).
\]
This is the expected dimension benchmark and deformation-theoretic lower bound for components through such curves. It specializes to the source's \(4d\) in \(\PP^3\) and \(5d+1-g\) in \(\PP^4\); the source says the latter is inadequate for fixed \(d\) and sufficiently large \(g\). Its original remark cites a better bound in the range \(d^3\le\lambda(n)g^2\), but does not define \(\lambda(n)\) or state that bound.

The supplied aid reports this component-construction criterion. Let \(2\le r\le n\), and let \(C\subset\PP^r\) be smooth, nondegenerate, and linearly normal, of degree \(d\) and genus \(g\), with
\[
 H^1(C,N_{C/\PP^r})=0.
\]
Linearly normal means that the restriction \(H^0(\PP^r,\mathcal O(1))\to H^0(C,\mathcal O_C(1))\) is an isomorphism. Moving the curve in its component and moving its linear span in \(\PP^n\) is reported to produce an actual ambient Hilbert component of dimension
\[
 h^0(C,N_{C/\PP^r})+(r+1)(n-r).
\]
When the original component has expected dimension, the ambient component's excess over \(E_n(d,g)\) is exactly
\[
 (n-r)(g+r-d).
\]
This is a family of exact benchmarks, not a universal improved lower bound.

For smooth degree-\(d\) plane curves with \(d\ge2\), \(g=(d-1)(d-2)/2\), the reported component in \(\PP^n\) has dimension
\[
 \binom{d+2}{2}-1+3(n-2).
\]
Its excess over \(E_n(d,g)\) and its normal obstruction-space dimension both equal
\[
 (n-2)\binom{d-2}{2}=h^1(C,N_{C/\PP^n}).
\]
Here a nonzero potential obstruction space is not a claim that actual deformations are obstructed; the aid calls it ineffective in this family. The restriction \(d\ge2\) retains the plane as the minimal span, excluding lines. Finally rational normal curves of degree \(r\) and genus zero in their span \(\PP^r\), including their ambient lifts, are reported to give an infinite family where expected dimension is sharp. The aid does not determine the requested sharp lower bound for all \((d,g,n)\).
\subsection{Short English statement}
The source asks for sharp lower bounds on dimensions of Hilbert-scheme components of curves. The supplied aid constructs exact component dimensions by varying a curve and its linear span, including plane-curve examples with excess dimension and rational normal examples where expected dimension is sharp.
\subsection{Sources}
\begin{itemize}
\item Local UnsolvedMath registry, permanent ID \(20000823\), source code \texttt{catalogue entry}; \texttt{aim-arithmetic-geometry-notes.json}, \texttt{notes[68]}, tag \texttt{problem}.
\item AIM workshop ``Components of Hilbert Schemes,'' 19--23 July 2010, Problem 22: \url{https://aimath.org/WWN/hilbertschemes/hilbertschemes.pdf}. Statement status \texttt{exact}; no PDF was accessed for this draft.
\item D. Chen, \emph{On the dimension of the Hilbert scheme of curves}, Math. Res. Lett. 16 (2009), no. 6, pp. 941--954, Theorem 1.3; cited in the original remark for the range \(d^3\le\lambda(n)g^2\).
\item Supplied research attempt 1: \texttt{partial\_result}, status \texttt{partially\_solved}, classification \texttt{PARTIAL-PROGRESS}, rigor label \texttt{proved}, novelty confidence low. Registry created and updated 14 August 2026. The candidate lifting criterion is not given a separate authoritative citation. This machine-generated aid is not peer reviewed and must be verified independently before citation.
\item Ulam AI / UnsolvedMath Contributors; CC BY 4.0.
\end{itemize}
\end{document}
